\documentclass[11pt]{article}
\usepackage{mystyle}

\title{Rosen --- \S 2.2: Set Operations\\ \large Selected Exercises}
\author{Alston}
\date{Semester 2, 2026}

\begin{document}
\maketitle

% ======================================================================
\begin{exercise}{15}
    Prove the second De Morgan law in Table 1 by showing that if $A$ and
    $B$ are sets, then $\overline{A \cup B} = \overline{A} \cap
    \overline{B}$
    \begin{enumerate}[label=(\alph*)]
        \item by showing each side is a subset of the other side.
        \item using a membership table.
    \end{enumerate}
\end{exercise}

% ======================================================================
\begin{exercise}{17}
    Show that if $A$, $B$, and $C$ are sets, then
    $\overline{A \cap B \cap C} = \overline{A} \cup \overline{B} \cup
    \overline{C}$
    \begin{enumerate}[label=(\alph*)]
        \item by showing each side is a subset of the other side.
        \item using a membership table.
    \end{enumerate}
\end{exercise}

% ======================================================================
\begin{exercise}{19}
    Show that if $A$ and $B$ are sets, then
    \begin{enumerate}[label=(\alph*)]
        \item $A - B = A \cap \overline{B}$.
        \item $(A \cap B) \cup (A \cap \overline{B}) = A$.
    \end{enumerate}
\end{exercise}

% ======================================================================
\begin{exercise}{31}
    Let $A$ and $B$ be subsets of a universal set $U$. Show that
    $A \subseteq B$ if and only if $\overline{B} \subseteq \overline{A}$.
\end{exercise}

% ======================================================================
\begin{exercise}{37}
    Show that if $A$ is a subset of a universal set $U$, then
    \begin{enumerate}[label=(\alph*)]
        \item $A \oplus A = \emptyset$.
        \item $A \oplus \emptyset = A$.
        \item $A \oplus U = \overline{A}$.
        \item $A \oplus \overline{A} = U$.
    \end{enumerate}
\end{exercise}

% ======================================================================
\begin{exercise}{41}
    Suppose that $A$, $B$, and $C$ are sets such that
    $A \oplus C = B \oplus C$. Must it be the case that $A = B$?
\end{exercise}

% ======================================================================
\begin{exercise}{51}
    Find $\bigcup_{i=1}^{\infty} A_i$ and $\bigcap_{i=1}^{\infty} A_i$ if
    for every positive integer $i$,
    \begin{enumerate}[label=(\alph*)]
        \item $A_i = \{-i, -i+1, \dots, -1, 0, 1, \dots, i-1, i\}$.
        \item $A_i = \{-i, i\}$.
        \item $A_i = [-i, i]$, that is, the set of real numbers $x$ with
        $-i \leq x \leq i$.
        \item $A_i = [i, \infty)$, that is, the set of real numbers $x$
        with $x \geq i$.
    \end{enumerate}
\end{exercise}

\end{document}